% Workshop 7 report skeleton. Limit: 4 pages in total.
% Put each screenshot in report/screenshots/ with the exact name used below.
% A missing image compiles as a grey box, so this file builds at any stage.
\documentclass[conference]{IEEEtran}

\usepackage[T1]{fontenc}
\usepackage{graphicx}
\usepackage{xcolor}
\usepackage{listings}
\usepackage{booktabs}
\usepackage{url}

\newcommand{\todo}[1]{\textcolor{red}{[TODO: #1]}}

% \shot{file}{caption}{label}: one-column figure.
\newcommand{\shot}[3]{%
  \begin{figure}[!t]
    \centering
    \IfFileExists{screenshots/#1}{%
      \includegraphics[width=\columnwidth]{screenshots/#1}%
    }{%
      \fcolorbox{gray}{gray!15}{\parbox[c][2.5cm][c]{0.92\columnwidth}{%
        \centering\footnotesize Missing: \texttt{screenshots/\detokenize{#1}}}}%
    }
    \caption{#2}
    \label{#3}
  \end{figure}%
}

\lstset{
  basicstyle=\ttfamily\scriptsize,
  columns=fullflexible,
  breaklines=true,
  frame=single,
  aboveskip=4pt,
  belowskip=4pt
}

\begin{document}

\title{Workshop 7: Data Interchange Using JSON\\
\large Schema-Validated and Schema-Less Services in a WhatsApp Farewell Scheduler}

\author{\IEEEauthorblockN{Ismael Paredes}
\IEEEauthorblockA{Distributed Systems, Semester II 2026 --- Prof. Francisco Hidrobo\\
Yachay Tech University\\
\todo{delivery date}}}

\maketitle

\begin{abstract}
\todo{3--4 sentences: what was built, the two services, the main measured result.}
\end{abstract}

\section{Objectives}
\todo{Exchange data between two components on different machines; define and enforce JSON Schemas for a request and its response; repeat without a schema; compare bytes, transfer latency and processing overhead.}

\section{Environment and Network Configuration}
\todo{Short description of the two machines and why traffic flows from Windows to WSL.}

\begin{table}[!t]
  \caption{Deployment}
  \label{tab:deployment}
  \centering
  \footnotesize
  \begin{tabular}{lll}
    \toprule
    & Machine A & Machine B \\
    \midrule
    Host     & Windows 11      & WSL2 Ubuntu \\
    IP       & \todo{IP}       & \todo{IP} \\
    Python   & \todo{version}  & \todo{version} \\
    Role     & Clients, benchmark & FastAPI services \\
    \bottomrule
  \end{tabular}
\end{table}

\shot{01_environments.png}{Hostname, IP address and Python version on both machines.}{fig:env}

\section{Services and Schemas}

\subsection{Activity 1: \texttt{POST /webhook} (with schema)}
\todo{Service description: the WhatsApp mock dispatcher delivers a customer message to the ingestion component, which assigns a partition and acknowledges. Where validation happens on each side.}

\begin{lstlisting}[caption={Excerpt of the request schema.},label={lst:schema}]
TODO: short excerpt of webhook_request.schema.json
\end{lstlisting}

\shot{02_schema_tests.png}{Schema test suite.}{fig:tests}
\shot{03_activity1_exchange.png}{Valid request and response on both machines.}{fig:a1}
\shot{04_validation_rejections.png}{Invalid data rejected by the client and by the server, and an invalid response detected.}{fig:reject}

\subsection{Activity 2: \texttt{POST /messages} (no schema)}
\todo{Service description: the scheduler dispatcher sends a message to the WhatsApp mock, which accepts it or answers 400/131047 outside the 24-hour window. Fields are read dynamically.}

\shot{05_activity2_exchange.png}{Schema-less request and response.}{fig:a2}
\shot{06_schemaless_bad_data.png}{Malformed data accepted by the schema-less service.}{fig:bad}

\section{Experiments and Results}
\todo{Method: cases A, A-control and B; number of requests; warm-up; keep-alive; how bytes, latency and overhead are measured.}

\begin{table}[!t]
  \caption{Results Summary}
  \label{tab:results}
  \centering
  \footnotesize
  \begin{tabular}{lrrr}
    \toprule
    Metric & A (schema) & A-control & B (no schema) \\
    \midrule
    Request body (B)        & -- & -- & -- \\
    Response body (B)       & -- & -- & -- \\
    Total with headers (B)  & -- & -- & -- \\
    RTT median (ms)         & -- & -- & -- \\
    RTT p95 (ms)            & -- & -- & -- \\
    Transfer latency (ms)   & -- & -- & -- \\
    Client overhead ($\mu$s) & -- & -- & -- \\
    Server overhead ($\mu$s) & -- & -- & -- \\
    \bottomrule
  \end{tabular}
\end{table}

\shot{07_benchmark_summary.png}{Benchmark summary output.}{fig:bench}
\shot{08_latency_overhead.png}{Latency and processing overhead by case.}{fig:chart}

\section{Problems and Solutions}
\todo{Real problems found during the workshop and how they were solved.}

\section{Discussion}
\todo{Cost of validation (A vs A-control); what the schema prevents (Fig.~\ref{fig:bad}); when each approach fits the project; limits: WSL2 shares the hardware, so latency is virtual-network latency.}

\section{Conclusions}
\todo{2--3 sentences.}

\section*{Source Code}
\url{https://github.com/mishellguanoc/DIST_WORKSHOPS/tree/main/workshop7}

\begin{thebibliography}{1}
\bibitem{jsonschema} JSON Schema, ``Draft 2020-12,'' \url{https://json-schema.org/draft/2020-12}.
\bibitem{meta} Meta, ``WhatsApp Business Platform Cloud API documentation.''
\end{thebibliography}

\end{document}
